
\documentclass[11pt,a4paper]{article}
\usepackage[utf8]{inputenc}
\usepackage[T1]{fontenc}
\usepackage[french]{babel}
\usepackage{lmodern}
\usepackage{microtype}
\usepackage[a4paper,margin=1.7cm,headheight=15pt]{geometry}
\usepackage{amsmath,amssymb,mathtools}
\usepackage{array,tabularx,multicol}
\usepackage{enumitem}
\usepackage[most]{tcolorbox}
\usepackage{xcolor}
\usepackage{tikz}
\usepackage{pgfplots}
\pgfplotsset{compat=1.18}
\usetikzlibrary{arrows.meta,calc,patterns}
\usepackage{fancyhdr}
\usepackage{fancyvrb}
\usepackage{listings}
\usepackage{hyperref}
\usepackage{pdfpages}

\definecolor{bleufonce}{RGB}{28,55,120}
\definecolor{bleuclair}{RGB}{238,243,255}
\definecolor{grisclair}{RGB}{246,246,246}
\definecolor{tableline}{RGB}{45,45,45}
\definecolor{softgray}{RGB}{246,247,249}
\definecolor{deepblue}{RGB}{25,62,115}
\definecolor{softblue}{RGB}{235,241,250}
\definecolor{accent}{RGB}{20,82,130}

\hypersetup{colorlinks=true,linkcolor=bleufonce,urlcolor=bleufonce,pdfauthor={},pdftitle={Parcours de révision -- Fonctions -- Énoncés}}
\setlength{\parindent}{0pt}
\setlength{\parskip}{0.28em}
\renewcommand{\arraystretch}{1.12}
\setlist[itemize]{leftmargin=1.25em,itemsep=0.15em,topsep=0.2em}
\setlist[enumerate,1]{label=\textbf{\arabic*.}, leftmargin=1.05cm, itemsep=0.35em, topsep=0.25em}
\setlist[enumerate,2]{label=\textbf{\alph*.}, leftmargin=0.85cm, itemsep=0.25em, topsep=0.2em}

\pagestyle{fancy}
\fancyhf{}
\lhead{}
\rhead{}
\cfoot{\thepage}

\newcommand{\e}{\mathrm e}
\newcommand{\dd}{\,\mathrm d}
\newcommand{\vect}[1]{\overrightarrow{#1}}
\newcommand{\origine}[1]{ {\scriptsize\textcolor{bleufonce}{(site #1)}}}
\newcounter{question}
\newcounter{reponse}

\tikzset{
  tabouter/.style={draw=tableline, line width=0.85pt},
  tabline/.style={draw=tableline, line width=0.55pt},
  forbidden/.style={draw=tableline, line width=0.95pt},
  vararrow/.style={-{Latex[length=2.7mm,width=2.0mm]}, line width=0.85pt, draw=accent},
  tabtext/.style={font=\small},
  tabmath/.style={font=\small},
  labelcell/.style={font=\small\bfseries, text=deepblue},
  pointvalue/.style={font=\small\bfseries},
  endvalue/.style={font=\small},
}

\newenvironment{blocquestions}[1]{%
  \begin{tcolorbox}[
    enhanced,
    breakable,
    colback=bleuclair,
    colframe=bleufonce,
    coltitle=white,
    colbacktitle=bleufonce,
    title={\bfseries #1},
    fonttitle=\bfseries,
    boxrule=0.6pt,
    arc=2mm,
    left=2mm,
    right=2mm,
    top=2mm,
    bottom=1.5mm,
    before skip=3mm,
    after skip=3mm, fontupper=\large
  ]
  \large
  \begin{enumerate}[leftmargin=*,label=\textbf{\arabic*.},itemsep=0.75em,topsep=0.5em]
  \setcounter{enumi}{\value{question}}
}{%
  \setcounter{question}{\value{enumi}}
  \end{enumerate}
  \end{tcolorbox}
}

\newenvironment{blocreponses}[1]{%
  \begin{tcolorbox}[
    enhanced,
    breakable,
    colback=grisclair,
    colframe=bleufonce,
    coltitle=white,
    colbacktitle=bleufonce,
    title={\bfseries #1},
    fonttitle=\bfseries,
    boxrule=0.6pt,
    arc=2mm,
    left=2mm,
    right=2mm,
    top=2mm,
    bottom=1.5mm,
    before skip=3mm,
    after skip=3mm, fontupper=\large
  ]
  \large
  \begin{enumerate}[leftmargin=*,label=\textbf{\arabic*.},itemsep=0.45em,topsep=0.5em]
  \setcounter{enumi}{\value{reponse}}
}{%
  \setcounter{reponse}{\value{enumi}}
  \end{enumerate}
  \end{tcolorbox}
}

\newtcolorbox{correctionbox}[1]{
  enhanced,
  breakable,
  colback=grisclair,
  colframe=black!55,
  boxrule=0.55pt,
  arc=2mm,
  left=2mm,right=2mm,top=1.5mm,bottom=1.5mm,
  title=\textbf{#1},
  coltitle=black,
  colbacktitle=black!12,
  before skip=2mm,
  after skip=2mm
}

\newcommand{\titreSujet}[2]{%
  \subsection*{#1}
  \addcontentsline{toc}{subsection}{#1}
  \textit{Référence : #2}\par\medskip\hrule\medskip
}

\lstset{language=Python,basicstyle=\ttfamily\small,numbers=left,numberstyle=\tiny,frame=single,showstringspaces=false,columns=fullflexible,keepspaces=true}
\sloppy

\begin{document}
\begin{center}
{\LARGE\bfseries Parcours de révision -- Fonctions}\\[1mm]

\end{center}
\vspace{1mm}
\tableofcontents
\clearpage

\phantomsection
\addcontentsline{toc}{section}{Partie 1 -- Récapitulatif de cours}
\includepdf[pages=-,pagecommand={},fitpaper=true]{parcours_revision_fonctions_carte_mentale_v15.pdf}

\clearpage
\section*{Partie 2 -- Questions flashs}
\addcontentsline{toc}{section}{Partie 2 -- Questions flashs}
\setcounter{question}{0}
\begin{blocquestions}{A. Vocabulaire essentiel}
  \item \origine{266} Qu'appelle-t-on une fonction affine ?
  \item \origine{267} Qu'appelle-t-on une fonction polynôme du second degré ?
  \item \origine{268} Qu'appelle-t-on une fonction rationnelle ?
  \item \origine{269} Qu'appelle-t-on le domaine de définition d'une fonction ?
  \item \origine{270} Qu'appelle-t-on une valeur interdite ?
  \item \origine{274} Qu'appelle-t-on une racine de l'équation $f(x)=0$ ?
  \item \origine{275} Qu'appelle-t-on une fonction dérivable sur un intervalle ?
  \item \origine{276} Qu'appelle-t-on le coefficient directeur d'une tangente ?
  \item \origine{277} Qu'appelle-t-on un extremum local ?
  \item \origine{280} Qu'appelle-t-on une fonction monotone sur un intervalle ?
\end{blocquestions}
\begin{blocquestions}{B. Limites, continuité et TVI}
  \item \origine{96} Que vaut $\displaystyle\lim_{x\to+\infty}(3x^2-5x+1)$ ?
  \item \origine{97} Que vaut $\displaystyle\lim_{x\to-\infty}\dfrac{2x-1}{x+3}$ ?
  \item \origine{98} Que vaut $\displaystyle\lim_{x\to2}\dfrac{x^2-4}{x-2}$ ?
  \item \origine{99} Qu'est-ce qu'une forme indéterminée ?
  \item \origine{103} Dans une fonction rationnelle, quel terme regarde-t-on d'abord pour la limite en $\pm\infty$ ?
  \item \origine{104} Qu'est-ce qu'une asymptote verticale ?
  \item \origine{105} Qu'est-ce qu'une asymptote horizontale à droite ?
  \item \origine{107} Que signifie qu'une fonction est continue en $a$ ?
  \item \origine{109} Quel théorème utilise-t-on pour montrer l'existence d'une solution de $f(x)=k$ ?
  \item \origine{110} Que faut-il souvent ajouter pour obtenir l'unicité d'une solution ?
\end{blocquestions}
\begin{blocquestions}{C. Exponentielle et logarithme}
  \item \origine{111} Comment note-t-on le logarithme népérien ?
  \item \origine{112} Quelle fonction est la réciproque de l'exponentielle ?
  \item \origine{113} Sur quel intervalle la fonction $\ln$ est-elle définie ?
  \item \origine{114} Que vaut $\ln(1)$ ?
  \item \origine{115} Que vaut $\ln(e)$ ?
  \item \origine{116} Que vaut $e^0$ ?
  \item \origine{117} Que vaut $\ln(e^x)$ ?
  \item \origine{118} Que vaut $e^{\ln x}$ pour $x>0$ ?
  \item \origine{119} Quelle propriété vérifie $\ln(ab)$ pour $a>0$ et $b>0$ ?
  \item \origine{120} Quelle propriété vérifie $\ln\left(\dfrac{a}{b}\right)$ pour $a>0$ et $b>0$ ?
\end{blocquestions}
\begin{blocquestions}{D. Dérivation, tangente, variations et convexité}
  \item \origine{125} Que vaut la dérivée de $e^x$ ?
  \item \origine{126} Que vaut la dérivée de $e^{ax+b}$ ?
  \item \origine{127} Que vaut la dérivée de $\ln(x)$ sur $]0,+\infty[$ ?
  \item \origine{128} Que vaut la dérivée de $\ln(ax+b)$ ?
  \item \origine{139} Que représente géométriquement $f'(a)$ ?
  \item \origine{140} Quelle est l'équation de la tangente à la courbe de $f$ au point d'abscisse $a$ ?
  \item \origine{141} Que signifie géométriquement $f'(a)=0$ ?
  \item \origine{142} Quelle information donne le signe de $f'(x)$ ?
  \item \origine{146} Qu'est-ce qu'une fonction convexe ?
  \item \origine{147} Quelle condition simple utilise-t-on souvent pour étudier la convexité ?
\end{blocquestions}
\begin{blocquestions}{E. Primitives, intégrales, équations différentielles et calculs flash}
  \item \origine{155} Qu'est-ce qu'une primitive de $f$ sur un intervalle ?
  \item \origine{157} Donner une primitive de $e^x$.
  \item \origine{160} Donner une primitive de $\dfrac1x$ sur $]0,+\infty[$.
  \item \origine{164} Quelle relation existe-t-il entre intégrale et primitive ?
  \item \origine{166} Qu'est-ce qu'une équation différentielle homogène du type $y'+ay=0$ ?
  \item \origine{167} Quelle est la forme générale des solutions de $y'+ay=0$ ?
  \item \origine{237} Dériver $x\mapsto 3x^2$.
\end{blocquestions}
\clearpage
\section*{Partie 3 -- Sujets de bac}
\addcontentsline{toc}{section}{Partie 3 -- Sujets de bac}
\titreSujet{Sujet 1 -- Exponentielle, équation différentielle, convexité et aire}{Asie, 5 septembre 2025, sujet 1, exercice 1}

Soit $f$ la fonction définie sur $\mathbb R$ par
\[
f(x)=xe^{-2x}.
\]
On admet que $f$ est deux fois dérivable sur $\mathbb R$ et on note $f'$ la dérivée de la fonction $f$.
On note $\mathcal C_f$ la courbe représentative de $f$ dans un repère orthonormé du plan.

Pour chacune des affirmations suivantes, préciser si elle est vraie ou fausse, puis justifier la réponse donnée. Toute réponse non argumentée ne sera pas prise en compte.

\noindent\textbf{Affirmation 1.}\quad Pour tout réel $x$, on a
\[
f'(x)=(-2x+1)e^{-2x}.
\]

\noindent\textbf{Affirmation 2.}\quad La fonction $f$ est une solution sur $\mathbb R$ de l'équation différentielle :
\[
y'+2y=e^{-2x}.
\]

\noindent\textbf{Affirmation 3.}\quad La fonction $f$ est convexe sur $]-\infty;1]$.

\noindent\textbf{Affirmation 4.}\quad L'équation $f(x)=-1$ admet une unique solution sur $\mathbb R$.

\noindent\textbf{Affirmation 5.}\quad L'aire du domaine délimité par la courbe $\mathcal C_f$, l'axe des abscisses et les droites d'équation $x=0$ et $x=1$ est égale à
\[
\dfrac14-\dfrac{3e^{-2}}4.
\]

\newpage
\titreSujet{Sujet 2 -- Logarithme, variations et convexité}{Polynésie, 18 juin 2025, sujet 2, exercice 2}

On considère la fonction $f$ définie sur l'intervalle $]2;+\infty[$ par
\[
f(x)=x\ln(x-2).
\]
Une partie de la courbe représentative $\mathcal C_f$ de la fonction $f$ est donnée ci-dessous.

\begin{center}
\begin{tikzpicture}
\begin{axis}[
width=11cm,height=6cm,
axis lines=middle,
xmin=1.8,xmax=8.5,
ymin=-15,ymax=15,
xtick={2,3,4,5,6,7,8},
ytick={-15,-10,-5,0,5,10,15},
grid=both,
samples=200,
clip=false,
]
\addplot[domain=2.05:8.5, thick] {x*ln(x-2)} node[pos=0.93, above] {$\mathcal C_f$};
\end{axis}
\end{tikzpicture}
\end{center}

\begin{enumerate}
  \item Conjecturer, à l'aide du graphique, le sens de variation de $f$, ses limites aux bornes de son ensemble de définition ainsi que les éventuelles asymptotes.
  \item Résoudre l'équation $f(x)=0$ sur $]2;+\infty[$.
  \item Calculer $\displaystyle\lim_{\substack{x\to 2\\ x>2}} f(x)$. Ce résultat confirme-t-il l'une des conjectures faites à la question 1 ?
  \item Démontrer que, pour tout $x$ appartenant à $]2;+\infty[$,
  \[
f'(x)=\ln(x-2)+\dfrac{x}{x-2}.
  \]
  \item On considère la fonction $g$ définie sur l'intervalle $]2;+\infty[$ par $g(x)=f'(x)$.
  \begin{enumerate}
    \item Démontrer que, pour tout $x$ appartenant à $]2;+\infty[$, on a :
    \[
g'(x)=\dfrac{x-4}{(x-2)^2}.
    \]
    \item On admet que $\displaystyle\lim_{\substack{x\to 2\\x>2}}g(x)=+\infty$ et que $\displaystyle\lim_{x\to+\infty}g(x)=+\infty$.

    En déduire le tableau des variations de la fonction $g$ sur $]2;+\infty[$. On fera apparaître la valeur exacte de l'extremum de la fonction $g$.
    \item En déduire que, pour tout $x$ appartenant à $]2;+\infty[$, $g(x)>0$.
    \item En déduire le sens de variation de la fonction $f$ sur $]2;+\infty[$.
  \end{enumerate}
  \item Étudier la convexité de la fonction $f$ sur $]2;+\infty[$ et préciser les coordonnées d'un éventuel point d'inflexion de la courbe représentative de la fonction $f$.
  \item Combien de valeurs de $x$ existe-t-il pour lesquelles la courbe représentative de $f$ admet une tangente de coefficient directeur égal à $3$ ?
\end{enumerate}

\newpage
\titreSujet{Sujet 3 -- Équation différentielle et étude d'une exponentielle}{Amérique du Sud, 21 novembre 2024, jour 1, exercice 1}

\textbf{Partie A}

On considère l'équation différentielle
\[
(E):\quad y'+\dfrac14 y=20e^{-\frac14 x},
\]
d'inconnue $y$, fonction définie et dérivable sur l'intervalle $[0;+\infty[$.

\begin{enumerate}
  \item Déterminer la valeur du réel $a$ tel que la fonction $g$ définie sur l'intervalle $[0;+\infty[$ par
  \[
g(x)=axe^{-\frac14 x}
  \]
  soit une solution particulière de l'équation différentielle $(E)$.
  \item On considère l'équation différentielle
  \[
  (E'):\quad y'+\dfrac14 y=0,
  \]
  d'inconnue $y$, fonction définie et dérivable sur l'intervalle $[0;+\infty[$. Déterminer les solutions de l'équation différentielle $(E')$.
  \item En déduire les solutions de l'équation différentielle $(E)$.
  \item Déterminer la solution $f$ de l'équation différentielle $(E)$ telle que $f(0)=8$.
\end{enumerate}

\textbf{Partie B}

On considère la fonction $f$ définie sur l'intervalle $[0;+\infty[$ par
\[
f(x)=(20x+8)e^{-\frac14 x}.
\]
On admet que la fonction $f$ est dérivable sur l'intervalle $[0;+\infty[$ et on note $f'$ sa fonction dérivée sur l'intervalle $[0;+\infty[$. De plus, on admet que
\[
\lim_{x\to+\infty} f(x)=0.
\]

\begin{enumerate}
  \item
  \begin{enumerate}
    \item Justifier que, pour tout réel $x$ positif,
    \[
f'(x)=(18-5x)e^{-\frac14 x}.
    \]
    \item En déduire le tableau de variations de la fonction $f$. On précisera la valeur exacte du maximum de la fonction $f$ sur l'intervalle $[0;+\infty[$.
  \end{enumerate}
  \item Dans cette question, on s'intéresse à l'équation $f(x)=8$.
  \begin{enumerate}
    \item Justifier que l'équation $f(x)=8$ admet une unique solution, notée $\alpha$, dans l'intervalle $[14;15]$.
    \item Recopier et compléter le tableau ci-dessous en faisant tourner étape par étape la fonction \texttt{solution\_equation} ci-dessous, écrite en langage Python.
  \end{enumerate}
\end{enumerate}

\begin{center}
\renewcommand{\arraystretch}{1.25}
\begin{tabular}{|c|c|c|c|c|c|}
\hline
$a$ & 14 & & & & \\ \hline
$b$ & 15 & & & & \\ \hline
$b-a$ & 1 & & & & \\ \hline
$m$ & 14,5 & & & & \\ \hline
Condition $f(m)>8$ & FAUX & & & & \\ \hline
\end{tabular}
\end{center}

\begin{Verbatim}[frame=single, fontsize=\small]
from math import exp

def f(x):
    return (20*x + 8)*exp(-1/4*x)

def solution_equation():
    a, b = 14, 15
    while b - a > 0.1:
        m = (a+b)/2
        if f(m) > 8:
            a = m
        else:
            b = m
    return a, b
\end{Verbatim}

\noindent\textbf{2.c.}\quad Quel est l'objectif de la fonction \texttt{solution\_equation} dans le contexte de la question ?

\newpage
\titreSujet{Sujet 4 -- Logarithme, continuité et suite associée}{Centres étrangers, 12 juin 2025, sujet 1, exercice 3}

\textbf{Partie A}

On considère la fonction $f$ définie sur l'intervalle $]-1;+\infty[$ par
\[
f(x)=4\ln(x+1)-\dfrac{x^2}{25}.
\]
On admet que la fonction $f$ est dérivable sur l'intervalle $]-1;+\infty[$.

\begin{enumerate}
  \item Déterminer la limite de la fonction $f$ en $-1$.
  \item Montrer que, pour tout $x$ appartenant à l'intervalle $]-1;+\infty[$, on a :
  \[
f'(x)=\dfrac{100-2x-2x^2}{25(x+1)}.
  \]
  \item Étudier les variations de la fonction $f$ sur l'intervalle $]-1;+\infty[$ puis en déduire que la fonction $f$ est strictement croissante sur l'intervalle $[2;6{,}5]$.
  \item On considère $h$ la fonction définie sur l'intervalle $[2;6{,}5]$ par $h(x)=f(x)-x$.

  On donne ci-dessous le tableau de variations de la fonction $h$ :

  \begin{center}
  \begin{tikzpicture}[x=1cm,y=1cm]
    \def\L{2.25}
    \def\C{3.00}
    \def\W{11.25}
    \def\H{2.70}
    \def\Yx{1.90}
    \def\Xm{5.25}

    % Fonds
    \fill[softgray] (0,\Yx) rectangle (\W,\H);
    \fill[softblue] (0,0) rectangle (\L,\Yx);

    % Cadre et séparations
    \draw[tabouter] (0,0) rectangle (\W,\H);
    \draw[tabline] (0,\Yx) -- (\W,\Yx);
    \draw[tabline] (\L,0) -- (\L,\H);
    \draw[tabline] (\Xm,0) -- (\Xm,\Yx);

    % Libellés
    \node[labelcell] at ({\L/2},{(\Yx+\H)/2}) {$x$};
    \node[labelcell] at ({\L/2},{\Yx/2}) {$h(x)$};

    % Ligne x
    \node[pointvalue] at ({\L+0.35},{(\Yx+\H)/2}) {$2$};
    \node[pointvalue] at (\Xm,{(\Yx+\H)/2}) {$m\simeq 2{,}364$};
    \node[pointvalue] at ({\W-0.45},{(\Yx+\H)/2}) {$6{,}5$};

    % Valeurs et flèches
    \node[endvalue] at ({\L+0.45},0.35) {$h(2)$};
    \node[pointvalue] at (\Xm,1.45) {$M\simeq 2{,}265$};
    \node[endvalue] at ({\W-0.60},0.35) {$h(6{,}5)$};

    \draw[vararrow] ({\L+0.90},0.50) -- ({\Xm-0.70},1.35);
    \draw[vararrow] ({\Xm+0.70},1.35) -- ({\W-1.15},0.50);
  \end{tikzpicture}
  \end{center}
  Montrer que l'équation $h(x)=0$ admet une unique solution $\alpha\in[2;6{,}5]$.
  \item On considère le script suivant, écrit en langage Python :
\begin{Verbatim}[frame=single, fontsize=\small]
from math import *

def f(x):
    return 4*log(1+x) - (x**2)/25

def bornes(n):
    p = 1/10**n
    x = 6
    while f(x) - x > 0:
        x = x + p
    return (x-p, x)
\end{Verbatim}
On rappelle qu'en langage Python :
\begin{itemize}
  \item la commande \texttt{log(x)} renvoie la valeur $\ln(x)$ ;
  \item la commande \texttt{c**d} renvoie la valeur de $c^d$.
\end{itemize}
  \begin{enumerate}
    \item Donner les valeurs renvoyées par la commande \texttt{bornes(2)}. On donnera les valeurs arrondies au centième.
    \item Interpréter ces valeurs dans le contexte de l'exercice.
  \end{enumerate}
\end{enumerate}

\textbf{Partie B}

Dans cette partie, on pourra utiliser les résultats obtenus dans la partie A.

On considère la suite $(u_n)$ définie par $u_0=2$ et, pour tout entier naturel $n$,
\[
u_{n+1}=f(u_n).
\]

\begin{enumerate}
  \item Montrer par récurrence que, pour tout entier naturel $n$,
  \[
  2\leqslant u_n\leqslant u_{n+1}<6{,}5.
  \]
  \item En déduire que la suite $(u_n)$ converge vers une limite $\ell$.
  \item On rappelle que le réel $\alpha$, défini dans la partie A, est la solution de l'équation $h(x)=0$ sur l'intervalle $[2;6{,}5]$. Justifier que $\ell=\alpha$.
\end{enumerate}

\newpage
\titreSujet{Sujet 5 -- Trigonométrie, convexité et intégration}{Amérique du Nord, 22 mai 2025, sujet 2, exercice 4}

On désigne par $f$ la fonction définie sur l'intervalle $[0;\pi]$ par
\[
f(x)=e^x\sin(x).
\]
On note $\mathcal C_f$ la courbe représentative de $f$ dans un repère.

\textbf{Partie A}

\begin{enumerate}
  \item
  \begin{enumerate}
    \item Démontrer que, pour tout réel $x$ de l'intervalle $[0;\pi]$,
    \[
f'(x)=e^x[\sin(x)+\cos(x)].
    \]
    \item Justifier que la fonction $f$ est strictement croissante sur l'intervalle $\left[0;\dfrac{\pi}{2}\right]$.
  \end{enumerate}
  \item
  \begin{enumerate}
    \item Déterminer une équation de la tangente $T$ à la courbe $\mathcal C_f$ au point d'abscisse $0$.
    \item Démontrer que la fonction $f$ est convexe sur l'intervalle $\left[0;\dfrac{\pi}{2}\right]$.
    \item En déduire que, pour tout réel $x$ de l'intervalle $\left[0;\dfrac{\pi}{2}\right]$,
    \[
e^x\sin(x)\geqslant x.
    \]
  \end{enumerate}
  \item Justifier que le point d'abscisse $\dfrac{\pi}{2}$ de la courbe représentative de la fonction $f$ est un point d'inflexion.
\end{enumerate}

\textbf{Partie B}

On note
\[
I=\int_0^{\frac{\pi}{2}}e^x\sin(x)\,dx\qquad\text{et}\qquad J=\int_0^{\frac{\pi}{2}}e^x\cos(x)\,dx.
\]

\begin{enumerate}
  \item En intégrant par parties l'intégrale $I$ de deux manières différentes, établir les deux relations suivantes :
  \[
  I=1+J\qquad\text{et}\qquad I=e^{\frac{\pi}{2}}-J.
  \]
  \item En déduire que
  \[
  I=\dfrac{1+e^{\frac{\pi}{2}}}{2}.
  \]
  \item On note $g$ la fonction définie sur $\mathbb R$ par $g(x)=x$.

  Les courbes représentatives des fonctions $f$ et $g$ sont tracées dans le repère orthogonal ci-dessous sur l'intervalle $[0;\pi]$.
  Calculer la valeur exacte de l'aire du domaine hachuré situé entre les courbes $\mathcal C_f$ et $\mathcal C_g$ et les droites d'équation $x=0$ et $x=\dfrac{\pi}{2}$.
\end{enumerate}

\begin{center}
\begin{tikzpicture}
\begin{axis}[
width=11cm,height=6cm,
axis lines=middle,
xmin=0,xmax=3.25,
ymin=0,ymax=9,
xtick={0,1,2,3},
ytick={0,2,4,6,8},
grid=both,
samples=200,
]
\addplot[domain=0:3.14159, thick] {exp(x)*sin(deg(x))} node[pos=0.80, above] {$\mathcal C_f$};
\addplot[domain=0:3.14159, thick] {x} node[pos=0.90, below] {$\mathcal C_g$};
\addplot[domain=0:1.5708, fill=gray!25, draw=none] {exp(x)*sin(deg(x))} \closedcycle;
\addplot[domain=0:1.5708, fill=white, draw=none] {x} \closedcycle;
\end{axis}
\end{tikzpicture}
\end{center}

\newpage
\titreSujet{Sujet 6 -- Équation différentielle, vitesse et intégration}{Métropole, 18 juin 2025, sujet 2, exercice 4}

L'objet de cet exercice est l'étude de l'arrêt d'un chariot sur un manège, à partir du moment où il entre dans la zone de freinage en fin de parcours.

On note $t$ le temps écoulé, exprimé en seconde, à partir du moment où le chariot arrive sur la zone de freinage.
On modélise la distance parcourue par le chariot dans la zone de freinage, exprimée en mètre, en fonction de $t$, à l'aide d'une fonction notée $d$ définie sur $[0;+\infty[$.
On a ainsi $d(0)=0$.
Par ailleurs, on admet que cette fonction $d$ est dérivable sur son ensemble de définition. On note $d'$ sa fonction dérivée.

\textbf{Partie A}

Sur la figure ci-dessous, on a tracé dans un repère orthonormé :
\begin{itemize}
  \item la courbe représentative $\mathcal C_d$ de la fonction $d$ ;
  \item la tangente $T$ à la courbe $\mathcal C_d$ au point $A$ d'abscisse $4{,}7$ ;
  \item l'asymptote $\Delta$ à $\mathcal C_d$ en $+\infty$.
\end{itemize}

\begin{center}
\begin{tikzpicture}
\begin{axis}[
width=11cm,height=6cm,
axis lines=middle,
xmin=0,xmax=12,
ymin=0,ymax=25,
xtick={0,2,4,6,8,10,12},
ytick={0,5,10,15,20,25},
grid=both,
samples=200,
]
\addplot[domain=0:12, thick] {exp(-0.6*x)*(-5*x/3 - 205/9) + 205/9} node[pos=0.78, above] {$\mathcal C_d$};
\addplot[domain=0:12, dashed] {205/9} node[pos=0.95, above] {$\Delta$};
\addplot[domain=2.3:7.2, thick] { (exp(-0.6*4.7)*(-5*4.7/3 - 205/9) + 205/9) + ((12+4.7)*exp(-0.6*4.7))*(x-4.7)} node[pos=0.95, above] {$T$};
\addplot[only marks, mark=*] coordinates {(4.7,{exp(-0.6*4.7)*(-5*4.7/3 - 205/9) + 205/9})} node[above left] {$A$};
\end{axis}
\end{tikzpicture}
\end{center}

Dans cette partie, aucune justification n'est attendue. Avec la précision que permet le graphique, répondre aux questions ci-dessous.

D'après ce modèle :
\begin{enumerate}
  \item Au bout de combien de temps le chariot aura-t-il parcouru $15$ m dans la zone de freinage ?
  \item Quelle longueur minimale doit être prévue pour la zone de freinage ?
  \item Que vaut $d'(4{,}7)$ ? Interpréter ce résultat dans le contexte de l'exercice.
\end{enumerate}

\textbf{Partie B}

On rappelle que $t$ désigne le temps écoulé, en seconde, à partir du moment où le chariot arrive sur la zone de freinage.
On modélise la vitesse instantanée du chariot, en mètre par seconde, en fonction de $t$, par une fonction $v$ définie sur $[0;+\infty[$.

On admet que :
\begin{itemize}
  \item la fonction $v$ est dérivable sur son ensemble de définition, et on note $v'$ sa fonction dérivée ;
  \item la fonction $v$ est une solution de l'équation différentielle
  \[
  (E):\quad y'+0{,}6y=e^{-0{,}6t},
  \]
  où $y$ est une fonction inconnue et où $y'$ est la fonction dérivée de $y$.
\end{itemize}
On précise de plus que, lors de son arrivée sur la zone de freinage, la vitesse du chariot est égale à $12$ m.s$^{-1}$, c'est-à-dire $v(0)=12$.

\begin{enumerate}
  \item
  \begin{enumerate}
    \item On considère l'équation différentielle
    \[
    (E'):\quad y'+0{,}6y=0.
    \]
    Déterminer les solutions de l'équation différentielle $(E')$ sur $[0;+\infty[$.
    \item Soit $g$ la fonction définie sur $[0;+\infty[$ par
    \[
g(t)=te^{-0{,}6t}.
    \]
    Vérifier que la fonction $g$ est une solution de l'équation différentielle $(E)$.
    \item En déduire les solutions de l'équation différentielle $(E)$ sur $[0;+\infty[$.
    \item En déduire que, pour tout réel $t$ appartenant à l'intervalle $[0;+\infty[$, on a :
    \[
v(t)=(12+t)e^{-0{,}6t}.
    \]
  \end{enumerate}
  \item Dans cette question, on étudie la fonction $v$ sur $[0;+\infty[$.
  \begin{enumerate}
    \item Montrer que, pour tout réel $t\in[0;+\infty[$,
    \[
v'(t)=(-6{,}2-0{,}6t)e^{-0{,}6t}.
    \]
    \item En admettant que
    \[
v(t)=12e^{-0{,}6t}+\dfrac{1}{0{,}6}\times 0{,}6te^{-0{,}6t},
    \]
    déterminer la limite de $v$ en $+\infty$.
    \item Étudier le sens de variation de la fonction $v$ et dresser son tableau de variation complet. Justifier.
    \item Montrer que l'équation $v(t)=1$ admet une solution unique $\alpha$, dont on donnera une valeur approchée au dixième.
  \end{enumerate}
  \item Lorsque la vitesse du chariot est inférieure ou égale à $1$ mètre par seconde, un système mécanique se déclenche permettant son arrêt complet.

  Déterminer au bout de combien de temps ce système entre en action. Justifier.
\end{enumerate}

\textbf{Partie C}

On rappelle que, pour tout réel $t$ appartenant à l'intervalle $[0;+\infty[$,
\[
v(t)=(12+t)e^{-0{,}6t}.
\]
On admet que, pour tout réel $t$ dans l'intervalle $[0;+\infty[$,
\[
d(t)=\int_0^t v(x)\,dx.
\]

\begin{enumerate}
  \item À l'aide d'une intégration par parties, montrer que la distance parcourue par le chariot entre les instants $0$ et $t$ est donnée par :
  \[
d(t)=e^{-0{,}6t}\left(-\dfrac53t-\dfrac{205}{9}\right)+\dfrac{205}{9}.
  \]
  \item On rappelle que le dispositif d'arrêt se déclenche lorsque la vitesse du chariot est inférieure ou égale à $1$ mètre par seconde.

  Déterminer, selon ce modèle, une valeur approchée au centième de la distance parcourue par le chariot dans la zone de freinage avant le déclenchement de ce dispositif.
\end{enumerate}
\end{document}
